% La residencia capitular pertenece a la estructura maestra. Este archivo
% desarrolla su contenido y conserva una segunda etiqueta compatible.
\label{chap:realizacion-espectral-numero-particula}
\index{number!spectral realization}
\index{persistent!section particle}
\index{route!observable shadow}
\index{spectrum!physical realization}

The HMT closure produces numbers as complete genealogical sections.
Dimensional Mechanics begins when a section of that class is realized on
a route, transports a register, induces a spectral operator, and acquires a
physical chart. This step preserves the mathematical antecedent and changes the
codomain: number, state, route, spectrum, and particle are objects linked
by explicit maps and retain different types.

The formal hinge is expressed through
\begin{equation}
 \boldsymbol x
 \xmapsto{\operatorname{sh}}
 \gamma_{\boldsymbol x}
 \xmapsto{\mathcal L}
 \mathcal L_{\boldsymbol x}
 \xmapsto{\Sigma}
 \sigma_{\boldsymbol x}
 \xmapsto{\chi}
 \chi_{\boldsymbol x}
 \xmapsto{\mathcal T}
 \widehat K_{\boldsymbol x}
 \xmapsto{\mathfrak R_{\rm phys}}
 \mathfrak p_{\boldsymbol x}.
 \label{eq:bisagra-numero-espectro-particula}
\end{equation}
Each arrow possesses its own domain and preserves the origin key of the initial section.

\section{Physical correspondence of the APP--TRIT--TPK core}
\label{sec:reaplicacion-fisica-app-trit-tpk}

The opening of Dimensional Mechanics retraces the formal core with
a new codomain. This repetition preserves the HMT theorems and declares
the physical-realization maps that are added. Its causal order is
\[
 \boxed{\begin{gathered}
 \text{APP as a support of sheets}
 \longrightarrow \text{TRIT as temporal orientation}\\
 \downarrow\\[-1mm]
 \text{unitary cycle and torsional defect}\\
 \downarrow\\[-1mm]
 \text{dual angular coordinates}
 \longrightarrow \text{particle--route}
 \end{gathered}}
\]

\subsection{The Cayley torus as a support for physical leaves}

APP provides the two-dimensional torus
\[
 X_9=(\mathbb Z/9\mathbb Z)^2
\]
with its Cayley graph and two evaluations on the same vertices. At the formal level, \(\Sigma\) accumulates and \(\Pi\) folds. The physical realization assigns sheets and channels to those evaluations, preserving the provenance map that makes it possible to compare their responses. The torus preserves spatial periodicity; the path category preserves direction, reversibility, and transport memory.

\subsection{Euclidean gnomonic chart and curvature atlas}
\label{sec:carta-gnomonica-equivalencia-ads-ds}

The L-shaped gnomon completes the multiplicative kernel \(8\times8\) and separates the
folded interior from the publication chart. In
the null class, precisely the two readers that the quotient identifies
outside the boundary are distinguished,
\[
 n\longmapsto n\bmod 9,
 \qquad
 n\longmapsto\operatorname{dr}_9(n),
\]
so that the residue \(0\) and its nonadic publication \(9\) preserve
distinct coordinates of the same closure. The sector \(3,6,9\) makes
this transition between the digital root and the modular residue visible.

In the geometric realization, the gnomonic L carries the Euclidean chart
\(\mathscr H_0^{(L)}\). On it, the inertial and gravitational readers
descend to a common normalization,
\[
 \frac{\mathcal G_{L}}{\mathcal I_{L}}=1,
\]
which constitutes the residence of the local equivalence principle. Before
that projection, the enriched state separately preserves the areal
and volumetric readings. The Ambivalence Principle is then expressed by
the internal ratio
\[
 \boxed{
 \frac{\mathcal L_{\rm ar}(x)}{\mathcal L_{\rm vol}(x)}
 =\frac{\pi_{\rm HMT}}{\varphi_{\rm HMT}^{2}}.}
\]
The Euclidean equality and the ambivalent ratio belong to distinct levels:
the former is the local projection onto the L; the latter preserves the
areal--volumetric genealogy prior to that projection.

The metric realization is obtained via a declared family of maps.
For
\[
 \mathbb A_s:=\mathbb R[X]/(X^2+s),
 \qquad s\in\{-1,0,+1\},
\]
and a curvature scale \(\rho>0\), one fixes
\[
 \mathcal M_{s,\rho}:\mathbb A_s\dashrightarrow
 (M_s,g_{s,\rho}),
 \qquad
 g_{s,\rho}=\mathrm d r^2+S_{s,\rho}(r)^2\mathrm d\theta^2,
\]
with
\[
 S_{s,\rho}(r)=
 \begin{cases}
  \rho^{-1}\sinh(\rho r),&s=-1,\\
  r,&s=0,\\
  \rho^{-1}\sin(\rho r),&s=+1.
 \end{cases}
 \qquad
 K(M_s,g_{s,\rho})=s\rho^2.
\]
The radial domain is \(r\geq0\) for \(s=-1,0\); for \(s=+1\), the polar chart covers \(0\leq r\leq\pi/\rho\), with the usual polar regularity at its endpoints. The branch \(s=-1\) provides the hyperbolic spatial section of the chart of anti--de Sitter type; \(s=0\) realizes the Euclidean gnomonic chart; and \(s=+1\) provides the positive spatial geometry that is prolonged into the de Sitter foliations. The complete correspondence is collected in \cref{sec:ansatz-metrico-tres-curvaturas,sec:atlas-ads-euclideo-ds-hmt}.

The algebraic classification and the metric family are coordinated through
the declared realization \(\mathcal M_{s,\rho}\). This arrow constitutes the
physical-interface ansatz, whereas
\(K=s\rho^2\) is the exact calculation of the metric once \(s\) and
\(\rho\) are fixed. The quadratic relation fixes the algebraic class; the declared
realization transports its sign to curvature, and \(\rho\) fixes the geometric
scale. In the negative chart, restriction to the boundary coordinates the
Archimedean projection with exceptional incidence; in the null chart, the
inertial and gravitational readers descend to local equivalence; in the
positive chart, the foliations satisfy \(\ddot a/a=H^2>0\). This atlas
links negative-curvature holography, local Euclidean physics, and
accelerated cosmological evolution through declared domains and change
maps.

\subsection{Temporal Representation of the TRIT: (+1) Return and Memory, (0) Threshold Present, and (-1) Bidirectional Opening}

Once an order has been declared on the transitions, the temporal reader is
\[
\begin{aligned}
 +1&\longmapsto\text{return, retained memory and past},\\
 0&\longmapsto\text{evaluation threshold and present},\\
 -1&\longmapsto\text{bidirectional opening and future}.
\end{aligned}
\]
The past designates transported closure; the present, the evaluation
boundary; the future, the tree of compatible prolongations. This reading
is the temporal realization of the three quadratic algebras already constructed
in HMT and preserves the involution that exchanges opening and return.

\subsection{Minimal Torsion and Unitary Dynamics in \texorpdfstring{\(\mathbb C^{108}\)}{C108}}

The invariant called minimal torsion is
\[
 \Delta_4
 =
 \frac{\pi-e\log\pi}{270}
 \left(1+\frac{\pi}{729}\right).
\]
Its entries arise from the HMT sections of \(\pi\) and \(e\), from the channel
\(270\), and from the interior \(729\). It functions as a global memory of
refinement before electron specialization.

The reference temporal dynamics resides in
\(\mathcal H_{\rm clk}=\mathbb C^{108}\). With
\[
 \omega_0=\frac{2\pi}{108\,t_0},
 \qquad
 H_{\rm clk}
 =\hbar_{\rm int}\omega_0
 \sum_{k=0}^{107}k\,|\widetilde k\rangle
 \langle\widetilde k|,
\]
the propagator satisfies
\[
 U_{\rm step}
 =\exp\!\left(-\frac{\mathrm i}{\hbar_{\rm int}}
 H_{\rm clk}t_0\right),
 \qquad
 U_{\rm step}^{108}=I.
\]
The unit cycle is exact for the declared primitives. Its realization
as a physical time crystal adds the Hamiltonian or driven protocol,
the primitive subharmonic response, stability under perturbations,
persistence, and the protection mechanism developed in
\cref{chap:regimenes-cuadraticos-tiempo}. This residence preserves the
hierarchy and refers the proof there.

\subsection{Dual angular coordinates and spectral channels}

The sum--product incompatibility is transported to two angular coordinates
\(A\) and \(C^*\). In the Archimedean chart
\[
 x=\frac{\pi A}{180},
 \qquad
 y=\frac{\pi C^*}{180},
\]
the channels are the eigenvalues of the central element
\[
 T=e^{-xI-yR}=q_+P_++q_-P_-,
 \qquad
 q_-=e^{-x+y},
 \qquad
 q_+=e^{-x-y}.
\]
The chart is invertible:
\[
 x=-\frac12\log(q_+q_-),
 \qquad
 y=\frac12\log\frac{q_-}{q_+}.
\]
Therefore, \(A\) and \(C^*\) are, respectively, the symmetric and
antisymmetric coordinates of a single spectral band. The
\cref{chap:canales} contains its construction and its proofs; the
present hinge fixes its causal position before the particle--route.

\section{State, route, spectral form and physical realization}

\begin{definition}[Finite state and genealogical section]
\label{def:estado-finito-seccion-genealogica-md}
A finite state is an element \(x_n\in\mathcal S_n^{\rm surv}\). A genealogical section is a
family \(\boldsymbol x=(x_n)_n\) compatible at every depth. The state is a restriction
of the section; the section determines its complete system of restrictions.
\end{definition}

\begin{definition}[Observable route]
\label{def:ruta-observable-bisagra-md}
A chart secuencial of depth \(R\) defines the sombra
\[
 \gamma_{\boldsymbol x}^{[R]}
 :=\operatorname{sh}_R(\boldsymbol x)
 =(c_0\to c_1\to\cdots\to c_R).
\]
The route preserves position, sheet, orientation, carry, and memory according to the
chart's data list. It is an ordered presentation of the section.
\end{definition}

\begin{definition}[Spectral form]
\label{def:forma-espectral-numero-md}
The spectral form associated with to section is
\[
 \mathfrak S(\boldsymbol x)
 :=(\gamma_{\boldsymbol x},
    \mathcal L_{\boldsymbol x},
    \sigma_{\boldsymbol x},
    \chi_{\boldsymbol x},
    \widehat K_{\boldsymbol x}),
\]
where \(\mathcal L_{\boldsymbol x}\) is the route record,
\(\sigma_{\boldsymbol x}\) its integer signature, \(\chi_{\boldsymbol x}\) the
character constructed on the signature, and \(\widehat K_{\boldsymbol x}\) the
operator obtained from the tower hierarchy. The spectrum
\(\operatorname{Spec}\widehat K_{\boldsymbol x}\) is an internal output.
\end{definition}

\begin{definition}[Physical realization]
\label{def:realizacion-fisica-forma-espectral}
A physical realization of \(\mathfrak S(\boldsymbol x)\) consists of
\[
 \mathfrak R_{\rm phys}
 =(\mathcal H_{\rm phys},\rho_{\rm sym},\varsigma_{\bf u},\mathcal O),
\]
where \(\mathcal H_{\rm phys}\) is the state space, \(\rho_{\rm sym}\) the representation of
the symmetry group, \(\varsigma_{\bf u}\) the units chart, and \(\mathcal O\) the set
of observables. The realization intertwines route transports with the
dynamics of \(\mathcal H_{\rm phys}\) and publishes quantities with units.
\end{definition}

The spectral form is exact within the discrete system once its
readers have been fixed. Identification with a physical species belongs to the
realization map. This separation makes it possible to publish the operator, its spectrum, its
position, and its status, accompanied by the corresponding physical intertwiner.

\section{\(4\) grammars and their domains}

The passage from number to particle uses four grammars. The
first three belong to the mathematical core; the fourth belongs to the
chosen statistical representation.

\begin{center}
\begin{tabular}{@{}lll@{}}
\toprule
\textbf{Grammar} & \textbf{Operation} & \textbf{Function}\\
\midrule
accumulation & sheet sum of APP & composes additive contributions;\\
fold & sheet product of APP & composes scales and transpositions;\\
orientation & TRIT and route conjugation & distinguishes sense and anti-section;\\
occupation & projector of the representation & fixes statistical multiplicity.\\
\bottomrule
\end{tabular}
\end{center}

The fermionic occupation is carried out, for example, in a fixed CAR algebra:
\[
 \{a_i,a_j\}=0,
 \qquad
 \{a_i^\dagger,a_j^\dagger\}=0,
 \qquad
 \{a_i,a_j^\dagger\}=\delta_{ij}I.
\]
It follows from these premises.
\[
 (a_i^\dagger)^2=0,
 \qquad
 n_i:=a_i^\dagger a_i,
 \qquad
 n_i^2=n_i,
 \qquad
 \operatorname{Spec}(n_i)=\{0,1\}.
\]
TPK selects the grading and transports the route; the CAR representation
realizes exclusion. This distinction positively fixes the algebraic
provenance of the Pauli principle.

\section{The spectral realization functor}

Let \(\Omega_{\rm HMT}^{\rm surv}\) be the category whose objects are surviving
sections and whose morphisms are transformations compatible with the
restrictions. Let \(\mathbf{SpecRoute}\) be the category of
\((\gamma,\mathcal L,\sigma,\chi,\widehat K)\) bundles and intertwiners that
preserve the route keys.

\begin{theorem}[Spectral realization of a section]
\label{thm:functor-realizacion-espectral-numero}
The route, register, signature, character, and tower readers define, on their
common domain, a functor
\[
 \mathfrak S:
 \Omega_{\rm HMT}^{\rm surv}\longrightarrow\mathbf{SpecRoute}.
\]
For a compatible restriction of sections, the route is truncated, the record
is restricted, and subsequent readers commute with that restriction. Two
presentations related by a genealogical symmetry produce intertwined spectral
forms.
\end{theorem}

\begin{proof}
The route shadow is defined by sequential reading of the state and commutes
with truncation. The record concatenates typed local data, and its
restriction eliminates only the posterior segment. The signature and the character
are calculated by declared homomorphisms over that record. The
tower hierarchy prolongs the character by composition. Identity and
composition are preserved at every stage.
\end{proof}

The theorem makes the expression “the number assumes spectral form” precise:
it transports a complete section to a route-and-spectrum package that preserves
its genealogy.

\section{Persistence and definition of particle}

For a prefix \(s_N\in\mathcal S_N^{\rm surv}\), let
\[
 \operatorname{Ext}_M(s_N)
 :=\{s_M\in\mathcal S_M^{\rm surv}:\rho_{M,N}(s_M)=s_N\},
 \qquad
 \mathcal L(s_N)
 :=\sup\{M:\operatorname{Ext}_M(s_N)\neq\varnothing\}.
\]

\begin{definition}[Persistence]
\label{def:persistencia-bisagra-numero-particula}
A prefix is persistent when it belongs to a compatible global section;
in a finitely branching tree of prolongations, this condition is equivalent
to \(\mathcal L(s_N)=\infty\).
\end{definition}

\begin{definition}[Mathematical particle]
\label{def:particula-matematica-bisagra}
Let \(\gamma_{\boldsymbol x}\) be a persistent route and let
\(\mathcal B_{\gamma_{\boldsymbol x}}\) be a discrete bundle whose fibres
carry the internal degrees of the realization. A mathematical particle is
\[
 \mathfrak p=
 (\boldsymbol x,
  \gamma_{\boldsymbol x},
  \mathcal B_{\gamma_{\boldsymbol x}},
  s,n,U,
  \mathfrak S(\boldsymbol x),
  \mathfrak R_{\rm phys}),
\]
where \(s\) is a compatible section of the bundle, \(n\) its
occupation function, \(U\) the propagator intertwining the transports, and
\(\mathfrak R_{\rm phys}\) the declared physical realization.
\end{definition}

The definition contains three closures:
\begin{enumerate}
\item the number is the global closure of compatible prefixes;
\item the spectral form is the closure of the route reading under the record,
      signature, character, and towers;
\item the particle is the local closure traversed by a persistent section
      of a physical bundle.
\end{enumerate}
The relation is summarized as
\[
 \boxed{
 \text{number}=\text{global genealogical closure},
 \qquad
 \text{particle}=\text{persistent local realization of that closure}.}
\]
The verbal equality denotes a functorial dependence between types.

\section{Oriented band and anti-section}

Let \(\tau=(\tau_0,\ldots,\tau_{11})\in\{-1,0,+1\}^{12}\) be the oriented word of a route, and define
\begin{equation}
 C_M(\tau):=-\operatorname{rev}(\tau).
 \label{eq:conjugacion-material-bisagra}
\end{equation}
It holds \(C_M^2=I\). The band
\[
 \mathcal B(\tau):=\{\tau,C_M\tau\}
\]
brings together two orientations of the same genealogy. One section realizes the
particle, and the conjugate section realizes the antiparticle.

\begin{proposition}[Mass parity and orientation]
\label{prop:paridad-masa-orientacion-bisagra}
Let \(\mathfrak m(\tau)\) be a mass reading that depends on even invariants
of the band, and let \(q(\tau)\) be an odd oriented reading. Then
\[
 \mathfrak m(C_M\tau)=\mathfrak m(\tau),
 \qquad
 q(C_M\tau)=-q(\tau).
\]
\end{proposition}

\begin{proof}
The first identity expresses invariance under the band involution. The
second expresses the parity of an oriented form. The relation \(C_M^2=I\)
guarantees the compatibility of both readings.
\end{proof}

Particle and antiparticle thus share the mass spectrum and have
opposite orientations. The conclusion follows from the parity of the readers.

\section{Path propagation and finite path summation}

Let \(U\) be a propagator on a finite state space. For every
\(R\geq1\), matrix multiplication gives
\begin{equation}
 (U^R)_{fi}
 =\sum_{\substack{\gamma:i\to f\\|\gamma|=R}}
   \prod_{\ell\in\gamma}U_\ell.
 \label{eq:feynman-finito-bisagra}
\end{equation}

\begin{theorem}[Exact expansion by histories]
\label{thm:feynman-finito-bisagra}
The equation \eqref{eq:feynman-finito-bisagra} enumerates once each route of
length \(R\) between \(i\) and \(f\). If \(U\) is unitary, \(U^R\) preserves
the norm and defines a finite unitary evolution.
\end{theorem}

\begin{proof}
Expanding the matrix product yields all the intermediate indices
\((x_1,\ldots,x_{R-1})\). Each tuple determines a route, and each route determines
a tuple. Unitarity follows from
\((U^R)^\dagger U^R=I\).
\end{proof}

The wave amplitude is the coherent sum over the band of routes; the
particle is the persistent section whose propagation realizes that band. Wave
and state form complementary levels of the same realization.

\section{Finite horizon of prolongation and class of compatible extension extensions of the genealogical state}

\begin{definition}[Virtual finite route]
\label{def:ruta-virtual-finita-bisagra}
A prefix \(s_N\) is finitely virtual if there exists \(L<\infty\) such that
\[
 \operatorname{Ext}_L(s_N)\neq\varnothing,
 \qquad
 \operatorname{Ext}_{L+1}(s_N)=\varnothing.
\]
Its route transports record, signature, and amplitude to \(L\), and its
prolongation horizon ends at the boundary \(L+1\).
\end{definition}

In a finitely branching tree, the existence of prolongations to every
depth is equivalent, by Koenig's lemma, to the existence of a global
section. Persistence and virtuality are defined through the
geometry of prolongations. The virtual route occupies a typed intermediate
state and has a finite genealogical duration.

\section{Internal spectrum, realization, and species}

For a multisection \(F\) of routes, the natural object is the operator
\[
 \widehat B_F e_\Gamma
 =\kappa(\Gamma)e_\Gamma,
 \qquad
 \exp((\log R_{\rm act})\widehat B_F).
\]
Its spectrum is invariant under permutations of the fiber basis. The physical
chart determines which combination of eigenvalues corresponds to a scalar
or matrix observable. Three levels must be preserved:
\[
 \operatorname{Spec}\widehat B_F,
 \qquad
 \mathfrak R_{\rm phys}(\widehat B_F),
 \qquad
 \text{species identification}.
\]
The first is an internal output; the second is a realization; the third
is a key within the physical inventory. A complete scientific table
publishes all three and the morphism relating them.

\section{Physical-spectral realization functor and 1st electronic closure}

The hinge is fixed by the chain
\[
 \boxed{\begin{gathered}
 \text{genealogical number}
 \longrightarrow\text{observable route}
 \longrightarrow\text{record}
 \longrightarrow\text{signature}
 \\
 \text{character}
 \longrightarrow\text{towers}
 \longrightarrow\text{spectral operator}
 \\
 \text{physical realization}
 \longrightarrow\text{particle}
 \end{gathered}}
\]
The first physical closure specializes this chain to the electron section selected by the center of the atlas, with the current action line, bidirectional readers, and dimensional chart. This residence establishes the formal domain that makes it possible to understand why the electron realizes the theory of the number and why its generalization to the other species preserves route, spectrum, and genealogy.

\paragraph{Composition of dimensional readers and compatibility among their mathematical invariants}
The definition of persistent section and discrete bundle belongs to the
particle chapter; route, register, signature, and character belong to the
surviving-extension chapter; band, material conjugation, and virtuality belong
to the antisections chapter; fiber operators and their spectra belong to the
bidirectional law of masses. The four grammars and the local closure of the
path are integrated here with the CAR structure, persistence, character,
and physical realization.
